\section{相关工作}
\label{sec:related-work}

\subsection{学习模型估计代价与辅助优化决策}

传统数据库优化器依赖人工设计的代价模型、基数估计器和搜索启发式来完成访问路径选择、连接顺序优化和物理算子选择~\cite{selinger1979access,chaudhuri1998overview}。
这类方法具有良好的工程可控性，但在复杂分析型查询中容易受到基数估计误差、代价模型偏差和搜索预算限制的影响。
因此，近年来大量工作尝试使用学习模型估计查询代价、执行时间或计划优劣，并将其嵌入优化器决策流程。

一类代表性方法将学习模型用于物理计划选择或计划排序。
Bao 通过树卷积模型学习优化器提示的收益，并结合 Thompson Sampling 在不同提示配置之间进行探索~\cite{marcus2021bao}。
Lero 使用成对排序模型预测候选执行计划的相对优劣，从而在数据库优化器生成的计划集合中选择更可能高效的计划~\cite{zhu2023lero}。
HybridQO 进一步考虑学习模型的不确定性，在传统优化器和学习型优化决策之间进行自适应切换，以提升部署鲁棒性~\cite{yu2022hybridqo}。
这些方法表明，学习模型能够从历史执行数据中捕获传统代价模型难以表达的性能模式，并在计划选择阶段改进优化质量。

另一类工作将学习信号用于逻辑重写或规则搜索。
例如，LearnedRewrite 将规则序列选择建模为学习型搜索问题，通过估计不同重写路径的潜在收益来指导规则组合~\cite{zhou2023learnedrewrite}。
与物理计划排序相比，这类方法更接近 SQL 重写本身，因为它们试图在逻辑表达式层面改变查询形态。
然而，学习式代价估计通常需要大量训练样本或离线标注执行结果，其有效性依赖训练负载、数据库实例和数据分布的覆盖范围。
当查询模板、数据分布或模式结构发生变化时，模型估计误差可能被规则序列搜索进一步放大，导致系统选择低收益甚至负收益的变换。
此外，学习模型的决策过程往往难以解释，给数据库内核集成、错误诊断和回退策略设计带来额外复杂性。

PASR 与上述工作互补。
它不训练一个全局代价模型来直接预测候选 SQL 的最终收益，也不只是在优化器已生成的物理计划之间排序。
相反，PASR 将执行计划分析结果转化为结构化的算子级瓶颈证据，并将其作为规则检索和规则序列编排的上下文。
这种设计保留了规则重写的可解释边界，同时避免将系统效果完全建立在离线学习模型的泛化能力之上。

\subsection{大语言模型指导 SQL 重写}

随着大语言模型在代码生成和语义推理任务中的能力提升，近期工作开始利用 LLM 辅助数据库交互、SQL 生成、优化建议和查询重写。
相较于传统学习模型，LLM 能够直接处理 SQL 文本、模式描述、自然语言规则、执行计划片段和历史示例等异构上下文，因此特别适合用于识别复杂查询中的潜在重写机会。
现有 LLM 辅助 SQL 重写方法大体可分为三类。

第一类方法通过离线微调、偏好优化或强化学习使模型适配查询优化任务。
例如，部分 LLM-based optimizer 将数据库元数据、SQL 和优化反馈序列化为训练样本，使模型学习生成更优查询或优化建议；E3-Rewrite 采用阶段化训练过程提升重写质量和执行收益~\cite{xu2025e3rewrite}。
这类方法能够增强模型对特定任务和工作负载的适配能力，但通常需要构造高质量训练数据，并面临跨数据库实例、跨模式和跨负载泛化的挑战。

第二类方法依赖提示工程或检索增强，在不重新训练模型的情况下向 LLM 提供重写规则、示例和上下文信息。
GenRewrite、LITHE、R-Bot 和 LLMR2 等系统利用自然语言规则说明、少样本示例或候选规则集合指导模型生成重写结果~\cite{sun2025rbot,dharwada2025lithe,li2024llmr2,liu2024genrewrite}。
其中，LLMR2 等方法通过限制模型输出规则标识或规则选择结果，降低开放式生成带来的非等价风险。
这类方法减少了训练成本，并能够借助 LLM 的语义推理能力处理复杂 SQL 结构；但如果缺乏对执行计划瓶颈的结构化刻画，模型仍可能难以判断哪些重写最可能改善当前查询的真实性能。

第三类方法进一步引入多智能体协作或外部工具反馈。
例如，QUITE 通过有限状态机组织多个 LLM agent，并结合数据库反馈对候选重写进行验证和迭代~\cite{song2025quite}。
这类系统强调将 LLM 放入可执行、可反馈的优化闭环中，而不是一次性生成 SQL。
不过，现有方法往往仍以单条查询为主要优化边界，历史正负经验难以被结构化沉淀并迁移到后续相似查询；同时，若直接将冗长原始执行计划输入提示，关键瓶颈算子和可改写锚点容易被大量弱相关计划信息淹没。

PASR 面向这些局限作出三点设计取舍。
首先，它通过 PlanAnalyzer 将原始执行计划压缩为算子级瓶颈、代价集中区域和可改写锚点，使 LLM 接收的是可操作的结构化计划证据，而不是高噪声长文本。
其次，PASR 不依赖完全开放式 SQL 生成，而是通过规则库、优化组调度、语法检查和语义一致性检查，将 LLM 的作用限制在规则检索、适用性判断和序列编排之内。
最后，PASR 引入 DTAM 双轨自适应记忆机制，将执行反馈沉淀为参数化规则偏好和非参数化案例记忆，使系统能够在无需离线微调的条件下跨查询复用经验。
因此，PASR 的核心区别并非简单使用 LLM，而是将 LLM 推理约束在计划引导、规则约束和反馈复用共同构成的优化闭环中。
